\dossier{352}{Eutychius, « croyant au Christ »}
\trad{(Vers) Ici a été placé un cruel trépas : dans sa première jeunesse, le jeune homme, arraché (à la vie), a franchi le seuil de la mort ; Eutychius, sage, sage, pieux et bienveillant, [qui], croyant [au Ch]rist, possède les récompenses de la lumière. Eutychius gît ici, lui-même l’unique honneur des siens. — Déposé le quatrième jour avant les nones d’octobre (4 octobre), sous le consulat de Théodose Auguste pour la troisième fois et d’Eugène (393) ; […] vingt-neuf ans et trois mois.}
\comm{Épitaphe métrique (CLE 1346), dont la pierre est perdue. Le nom du Christ est lu à partir de \emph{-ristum} ; la restitution \emph{[qui Ch]ristum}, contrainte par la finale, le sens et le mètre, est très probable, mais elle reste une restitution (chapitre~\ref{ch:typologie}). Le participe \emph{credens} construit avec l’accusatif exprime l’adhésion personnelle ; la récompense est « la lumière ». La date (\emph{Theodosio Aug. III et Eugenio}) est celle d’un consulat reconnu à Rome sous l’usurpation d’Eugène. Les graphies \emph{iubenta}, \emph{iubenis} et \emph{Eutichius}, normalisées par l’EDCS, sont conservées par l’EDB ; les deux bases corrigent \emph{benignus}, gravé avec un C.}

\dossier{378}{Une généalogie cléricale transmise par un seul témoin}
\trad{[Limenius] vécut […] ans ; […] Adeodata (?) […] ; […] il est tenu pour sage (?) […] ; […] doux pour moi d’abord […] ; […] il repose […] ; […] mois, à son épouse […] ; déposé le onzième jour avant les calendes d’octobre (21 septembre), [sous le consulat d’Olybrius et de Pro]binus (395). [Sa fille Vict]oria, déposée aux ides de janvier (13 janvier), [avec son mari …]orius. Telle est la lignée de la fille : Limenius, son père, prêtre ; petit-fils de Crescentius, petit-fils de […], diacre de la cinquième région ; petit-fils de […], arrière-petit-fils de l’évêque Marcellus, petit-fils de Crescentius, \emph{matrona} d’un prêtre ; arrière-arrière-petit-fils de Basilius ; belle-mère de l’évêque […], évêque Ursinus d’Albano.}
\comm{Ce texte, transmis par l’EDCS sous les références ICVR III 8161 et 8345, pose un problème de tradition (chapitre~\ref{ch:evenements}). Les fiches de l’EDB qui se rapportent au monument (ICVR III 8345, 8274c et 8387v) ne contiennent ni la formule consulaire ni la généalogie ; l’EDB date l’ensemble « 390–424 ». L’année 395 et toute la généalogie reposent donc sur un seul témoin, dont plusieurs termes sont restitués. La traduction suit l’ordre du texte transmis sans prétendre en reconstituer la syntaxe : l’accord des termes de parenté (\emph{nepos}, \emph{pronepos}, \emph{abnepos}, \emph{socrus}) et leur rattachement à la fille ou à d’autres personnes restent incertains. Si la lecture est confirmée, le texte atteste une famille où se succèdent prêtres, diacres et évêques, et le mariage des clercs. Contrôle de l’ICVR III 8345 indispensable.}

\dossier{382}{Panthagathè (?), mère de trois enfants}
\trad{(Vers acrostiches) [Avec la pal]me, elle possède les royaumes célestes […] ; [mais] d’une chasteté entière […] ; naguère, une si grande fidélité en toi […] ; une si grande pudeur, l’esprit […] ; cette âme reconnaissante […] ; […] l’enfant né avant terme, combien […] ; épouse généreuse pour moi et […] ; dorées sont les (demeures ?) célestes que toi, au ciel […] ; mère de trois enfants, [elle repose dans le sein de Dieu] ; son nom, [les initiales] des vers [le déclarent]. Elle vécut vingt-cinq ans ; [déposée] le dix-septième jour avant les calendes de septembre (16 août), sous le consulat d’Olybrius et de Probinus (395).}
\comm{Le poème est conservé sur sa moitié gauche ; les initiales des vers formaient le nom de la défunte, qui ne peut être établi sur les lettres conservées. L’EDCS ajoute, sur une autre face ou une autre partie du support, \emph{[Pa]nthagate}, absent de l’EDB. Le ciel (\emph{celestia regna}, \emph{caelestis}, \emph{caelo}) et la fidélité (\emph{fides}) sont lus ; Dieu n’apparaît que dans la restitution \emph{in sin[u dei pausat]}, commune aux deux bases (chapitre~\ref{ch:typologie}). Le vers \emph{acerbus natus} suggère une mort liée à une naissance prématurée, sans que la lacune permette de l’affirmer.}

\dossier{387}{Un achat sous l’épiscopat de Sirice ?}
\trad{[…] déposée le dix […], [sous le consulat de notre seigneur Honorius Auguste pour la troisième fois et] de notre seigneur Flavius Arcadius Auguste pour la quatrième fois (396) ; […] Eros, Leo et Setema [ont acheté] un emplacement [près du seuil de Pierre et de] Paul, sous l’épiscopat de Si[rice, aux prêtres Ursus] et Proclinus.}
\comm{Dossier examiné au chapitre~\ref{ch:sepultures}. Sont lus \emph{locu(m)}, \emph{Pauli sub episcopatu Si-} et \emph{et Proclino}. Le verbe d’achat, la mention de Pierre, la qualité de prêtres des vendeurs et le nom d’Ursus sont des restitutions. La datation épiscopale est vraisemblable (Sirice est évêque de 384 à 399) ; la localisation près de la tombe de Paul et l’intervention du clergé dans la vente ne le sont qu’au prix de ces restitutions. Le texte n’est pas apparié à l’EDB.}

\dossier{388}{Quodvultdeus}
\trad{Ici repose Quodvultdeus, homme d’honorable mémoire, qui vécut cinquante-sept ans ; déposé en paix le cinquième jour avant les ides d’octobre (11 octobre), sous le consulat — (EDB seule :) de nos seigneurs Arcadius Auguste pour la quatrième fois et Honorius Auguste pour la troisième fois (396).}
\comm{\emph{Quodvultdeus}, « ce que Dieu veut », est un anthroponyme : il contient le nom de Dieu sans constituer une invocation (chapitre~\ref{ch:noms}). L’EDCS s’arrête à \emph{conss(ulibus)} ; l’EDB transmet les noms consulaires, qui fondent l’année 396 (chapitre~\ref{ch:fondement}). La plaque (71 × 145 cm) est conservée au Lapidario Cristiano ; l’EDB décrit une feuille.}

\dossier{393}{Deux fragments d’une date consulaire}
\trad{[Sous le consulat d’Arcadius pour la] quatrième fois [et d’]H[onorius pour la troisième fois, Augustes] (396).}
\comm{L’EDCS publie sous deux numéros (anciens \nnums{393} et 394) des fragments que l’EDB réunit sous une fiche unique (ICVR IV 11126 a-b ; 16 × 13 × 2 cm). L’EDB ne lit que \emph{…IIII H…}, \emph{…ADIO…} et \emph{…o III Aug(usto)…}, compatibles avec le consulat d’Arcadius IV et Honorius III. Les deux fragments sont réunis sous ce numéro, et un seul texte daté de 396 est compté (chapitre~\ref{ch:corpus}). L’année repose sur une restitution guidée par les itérations (niveau R).}

\dossier{400}{Severianus : « dans la lumière du Seigneur »}
\trad{Doux et innocent, ici dort Severianus dans le sommeil de la paix ; il vécut environ cinquante ans ; son esprit a été accueilli dans la lumière du Seigneur le [quantième anomal] jour avant les calendes de mars, sous le consulat de nos seigneurs Flavius Caesarius et Nonius Atticus, \emph{viri clarissimi} (397) ; cet emplacement, Viricunda, son épouse, l’a acheté — Pascasius, \emph{ep(is)c(opus)} (?).}
\comm{Texte mixte (épitaphe et acte d’achat ; chapitre~\ref{ch:sepultures}). L’énoncé religieux est lu dans les deux bases : le défunt « dort dans le sommeil de la paix » et son esprit a été « accueilli dans la lumière du Seigneur ». L’EDB décrit, dans le texte même, un chrisme après le nom du défunt et une croix monogrammatique après \emph{susceptus est}, et, à la fin, un oiseau portant un rameau d’olivier et une croix monogrammatique dans une couronne. Le quantième \emph{VIIIIX} est anomal ; l’EDB reprend la note de Diehl, qui propose \emph{prid. Id. Febr.} et observe qu’aucun évêque de Rome ne porte le nom de Pascasius. L’EDB ne développe pas \emph{EPC}. La fonction de ce Pascasius dans l’acquisition reste indéterminée, et la traduction laisse la construction ouverte. L’âge, \emph{plus minus L}, ne comporte pas de jours : les « dix-sept jours » de la table des personnes provenaient de la date (chapitre~\ref{ch:personnes}).}

\dossier{405}{Un catéchumène}
\trad{À Bonifa(tius ?), qui repose bien, catéchumène, qui vécut environ soixante ans ; déposé le quatrième jour avant les calendes de juillet (28 juin), sous le consulat de Caesarius et d’Atticus (397).}
\comm{Seul texte du corpus qui nomme un catéchumène (chapitre~\ref{ch:foi_eglise}). Le défunt, mort à environ soixante ans, n’est pas dit baptisé ; le statut est gravé comme une qualification. Le nom est transmis sous la forme \emph{Bonifae}, que l’on peut comprendre comme une abréviation ou une graphie de \emph{Bonifatius}. L’EDB transcrit une partie effacée (\emph{〚- - -IV- - -〛}) et une addition (\emph{《d(epositus)》}) qui pourraient signaler une correction sur la pierre ; la plaque est publiée avec une reproduction dans un catalogue de 2000.}

\dossier{411}{Ursa : un prix d’emplacement}
\trad{Déposée en paix la jeune fille du nom d’Ursa, âgée de vingt-deux ans, sous le consulat de Flavius Caesarius et de Maximus (?), \emph{viri clarissimi} ; emplacement acheté (pour) trois (?) et onze sous.}
\comm{Texte mixte. L’âge est \emph{annorum viginti duo}, vingt-deux ans, et non vingt comme le portait la table des personnes ; \emph{puella} désigne ici une jeune femme (chapitre~\ref{ch:personnes}). Le collège transmis, \emph{Fl(avi) Caesari et Maximi}, ne correspond pas aux fastes de 397 (Caesarius et Atticus) : l’année est conservée sous réserve. Le prix (\emph{(t)res et s(olidis) XI} dans l’EDCS, \emph{⟨t⟩ris et s(olidis) (undecim)} dans l’EDB) est le seul du corpus ; sa construction reste incertaine (chapitre~\ref{ch:sepultures}).}

\dossier{433}{Un petit enfant parmi les astres}
\trad{(Vers) Ici gît, éteint au premier [seuil de la vie], un enfant très aimé [pendant un temps] bien court ; et je ne pense pas [qu’il soit permis d’accabler de] larmes [ce bienheureux], qui, dépouillé des liens du corps, [se réjouit parmi les astres], et n’a pas ressenti les [contagions] mauvaises [du sens] terrestre ; car après [avoir ajouté] neuf mois [à une année], il a accompli six jours de vie (de plus) [et s’en est allé vers les astres]. — Déposé sous le consulat de Lucius Mallius Theodorus, \emph{vir clarissimus} (399), la veille des nones […] ; il vécut un an, neuf mois [et six jours].}
\comm{Les vers sont conservés sur leur moitié gauche, et l’EDB souligne de nombreuses lettres de cette moitié, lues par les éditeurs anciens et aujourd’hui perdues. Les mots \emph{astris} et \emph{astra}, qui portent l’image de l’ascension, appartiennent entièrement aux fins de vers restituées (chapitre~\ref{ch:mort_esperance}). Ce qui est lu (\emph{corporis exutus vinclis}, \emph{nec mala terreni sensit}) suffit cependant à situer le poème dans le registre de la consolation, qui présente la mort de l’enfant comme une libération. L’EDB lit \emph{angusto} (« bref ») là où l’EDCS lit \emph{augusto}. L’année repose sur le consulat de Mallius Theodorus, lu ; le prénom \emph{Lucius} gravé est inhabituel pour ce consul.}

\dossier{447}{Une imprécation grecque}
\trad{(Grec ; partie transmise par l’EDB seule :) [Ici] repose Eu[…], [citoyen de Théo]dosiopolis, [ayant vécu] un peu plus de […] ans, [il s’est endormi] le quinzième jour avant les calendes de septembre (18 août), [sous le consulat de Vincentius] et de Fravitta, [clarissimes] (401). — (Partie commune aux deux bases :) Je vous adjure par [Dieu le Tout-]puissant : si quelqu’un [ose déranger] ses os [ou …], qu’il soit passible [du jugement / du châtiment] […] au [redoutable] jour [de la résurrection (EDCS) / du jugement (EDB)].}
\comm{Dossier examiné aux chapitres~\ref{ch:grec} et~\ref{ch:mort_esperance}. L’EDCS ne transmet que les cinq dernières lignes ; l’EDB en donne dix, qui contiennent le nom, l’origine et la date, et suit une autre édition. Dans la formule d’adjuration, sont lus \emph{ὁρκίζομε δὲ ὑμᾶς}, \emph{-ντοκράτορα}, \emph{εἴ τις το-}, \emph{τὰ ὀστέα} (EDB : \emph{ᾡτέα}) \emph{αὐτοῦ}, \emph{ἔνοχος τῇ} et \emph{ἡμαίραν τῆς}. Les deux bases restituent différemment la sanction : \emph{τῇ [κρίσει] … τῆς [ἀναστάσεως(?)]} dans l’EDCS, \emph{τῇ [κολάσει τὴν φοβερ]ὰν ἡμαίραν τῆς [κρίσεως]} dans l’EDB. Le texte atteste une menace eschatologique ; il n’atteste ni la résurrection ni le jugement par leur nom. L’expression \emph{ἔνοχος … τῇ κρίσει}, si elle est retenue, rappelle Mt 5, 21–22.}

\dossier{458}{Praetiosa, \emph{ancilla dei et Christi}}
\trad{La veille des ides de juin (12 juin) a reposé Praetiosa, jeune fille, vierge, de douze ans seulement, servante de Dieu et du Christ, sous le consulat de Flavius Vincentius et de Fravitta, \emph{viri clarissimi} (401).}
\comm{La formule \emph{ancilla dei et Chr(ist)i} associe Dieu et le Christ dans une qualification de dévotion ; elle ne contient ni « Père » ni « Fils » (chapitre~\ref{ch:dieu_christ}). L’EDB transcrit \emph{Chr(ist)i} : le I final est gravé, et il s’agit d’une abréviation en lettres, non d’un monogramme. L’association de \emph{virgo} et d’\emph{ancilla dei} pour une jeune fille de douze ans peut signaler une consécration précoce ou un éloge de l’innocence ; le texte ne permet pas de choisir. Le nom est \emph{Praetiosa} dans l’EDB, \emph{Praetioso} dans l’EDCS. La plaque provient, selon l’EDB, de la crypte de saint Corneille, dans le complexe de Calliste.}

\dossier{470}{Felicitas, « appelée d’un surnom trompeur »}
\trad{(Vers) Ce qu’offrent un doux enfant, ce qu’offrent de chers gages, ce tombeau le contient, qui retient de petits membres. Elle a causé une douleur sans fin, Felicitas, enfermée ici, malheureuse, appelée d’un surnom trompeur ; car, à elle que le premier jour avait mise au monde, (furent données) quatre années, et, comme elle commençait malheureusement la cinquième, un jour noir l’emporta et la plongea dans une mort prématurée. Sa mère et son père composent ces vers pour le tombeau, afin qu’en même temps, pour qui les lit, l’image revienne au cœur, et que les yeux, souvent séchés de pleurs, se mouillent de nouveau ; ainsi l’amour se soigne, et les Mânes n’ont que faire des vers. — Déposée en paix aux calendes de juillet (1\textsuperscript{er} juillet), sous le consulat de notre seigneur Théodose Auguste et de Flavius Rumoridus, \emph{vir clarissimus} (403).}
\comm{Épitaphe métrique (CLE 682) gravée sur un sarcophage, en deux parties séparées, selon l’EDB, par l’image d’une jeune orante. Le poème est entièrement classique dans sa langue et ses motifs : le jeu sur le nom (\emph{Felicitas} / \emph{infelix}), le \emph{funus acerbum}, la fonction consolatrice des vers, et même les Mânes, qui « n’ont que faire des vers ». La partie en prose se termine par \emph{deposita in pace}. Le texte illustre la coexistence, dans un même monument chrétien, du langage poétique traditionnel et du formulaire chrétien (chapitre~\ref{ch:mort_esperance}).}

\dossier{474}{Quintilianus, \emph{confirmans trinitatem}}
\trad{Quintilianus, homme de Dieu, […], affermissant (ou professant) la Trinité, […], aimant la chasteté, […], rejetant le monde, […], repose avec […], qui vécut environ dix […] ans ; sous le consulat de Théodose Auguste et de R[umoridus, \emph{vir clarissimus}] (403).}
\comm{Dossier examiné au chapitre~\ref{ch:esprit_trinite}. Le texte, identique dans les deux bases, est conservé sur sa moitié gauche. Les trois participes forment un éloge ascétique ; \emph{confirmans trinitatem} qualifie la foi du défunt, sans nommer les trois personnes. L’âge commence par \emph{X} et se poursuit dans la lacune : il est d’au moins dix ans. L’année repose sur \emph{Theodosio Aug(usto)}, lu, et sur l’initiale \emph{R}. La pierre, aujourd’hui perdue, provient selon l’EDB du cimetière à ciel ouvert proche de la basilique Saint-Pierre.}

\dossier{475}{\emph{Dedit deus, deus tulit}}
\trad{[…] Dieu a donné, Dieu a repris ; […] des calendes de novembre, sous le consulat de Théodose et de Rumoridus, \emph{viri clarissimi} (403).}
\comm{Citation libre de Jb 1, 21 (Vulgate : \emph{Dominus dedit, Dominus abstulit}), employée comme consolation (chapitre~\ref{ch:dieu_christ}). La forme \emph{deus … tulit} peut dépendre d’une version latine antérieure à la Vulgate ou d’une citation de mémoire. Le texte est un fragment ; l’EDB le rattache au cimetière d’Octavilla, près de Saint-Pancrace.}
